\documentclass{article}

\usepackage[final]{neurips_2024}

\usepackage[utf8]{inputenc}
\usepackage[T1]{fontenc}
\usepackage{hyperref}
\usepackage{url}
\usepackage{booktabs}
\usepackage{amsfonts}
\usepackage{nicefrac}
\usepackage{microtype}
\usepackage{xcolor}
\usepackage{graphicx}
\usepackage{amsmath}
\usepackage{amssymb}
\usepackage{subfigure}
\usepackage{algorithm}
\usepackage{algpseudocode}
\usepackage[capitalize,noabbrev]{cleveref}

\title{Robust Motion Imitation under Force Perturbations: \\ A Study of AMP and ADD on Humanoid Robots}

\author{%
  Anirudh Shrihari \\
  Carnegie Mellon University\\
  \texttt{ashriha2@andrew.cmu.edu} \\
  \And
  Prajwal Gurunath \\
  Carnegie Mellon University\\
  \texttt{pgurunat@andrew.cmu.edu} \\
  \And
  Manyung Emma Hon \\
  Carnegie Mellon University\\
  \texttt{mehon@andrew.cmu.edu} \\
}

\begin{document}

\maketitle

\begin{abstract}
Motion imitation learning methods for humanoid robots typically showcase impressive feats like parkour, dancing, and running, but these skills are often demonstrated in ideal conditions without real-world disturbances. We investigate the robustness of two adversarial imitation learning methods—Adversarial Motion Priors (AMP) and Adversarial Differential Discriminator (ADD)—when subjected to external force perturbations representing practical tasks such as carrying objects. Our goal is to train policies that maintain natural human-like gaits while handling box-carrying forces. We implemented a comprehensive evaluation framework including: (1) custom reference motions for the Unitree G1 humanoid simulating box-carrying poses, (2) curriculum-based force perturbation training with both constant and random forces, (3) force decay mechanisms for realistic dynamics, (4) multi-directional force configurations [10N, 10N, 30N] to mitigate overcompensation, and (5) steering tasks for both AMP and ADD to evaluate directional control under load. Our experiments show that ADD learns faster and achieves longer rollouts but tends to overcompensate with unnatural poses under single-axis forces, while AMP maintains better motion quality but struggles with tracking accuracy. Multi-directional force training significantly reduces overcompensation behaviors. We provide insights into the trade-offs between these methods and demonstrate effective strategies for robust real-world deployment.
\end{abstract}

\section{Introduction}

Physics-based character animation and motion imitation have made remarkable progress with deep reinforcement learning~\citep{peng2018deepmimic,peng2021amp,zhang2025add}. Modern approaches can replicate complex human motions including acrobatics, dance, and athletic movements with high fidelity. However, these methods are typically evaluated in pristine simulation environments without the perturbations and disturbances common in real-world scenarios.

Real-world robotic applications demand more than motion reproduction—they require robustness to external forces, payload handling, and environmental variations. A humanoid robot must maintain stability and natural movement while carrying objects, responding to unexpected pushes, or operating on varied terrain. This gap between impressive demonstrations and practical utility motivates our investigation.

\textbf{Problem Formulation.} We focus on a specific but representative challenge: training humanoid robots to walk naturally while carrying a box. This task combines several difficulties: (1) continuous external forces on the end-effectors, (2) need to maintain human-like gait patterns from reference motions, and (3) robust recovery from varying force magnitudes. We evaluate two state-of-the-art adversarial imitation learning methods: AMP~\citep{peng2021amp} and ADD~\citep{zhang2025add}.

\textbf{Contributions.} Our work makes the following contributions:
\begin{itemize}
\item A systematic evaluation of AMP and ADD under external force perturbations, including constant forces (box carrying) and random impulses
\item Custom dataset generation for box-carrying reference motions on the Unitree G1 humanoid
\item Curriculum learning framework for progressive force adaptation with configurable decay dynamics
\item Implementation of random force perturbations and steering tasks extending the MimicKit framework
\item Empirical analysis revealing distinct failure modes and trade-offs: ADD's speed vs. naturalness, AMP's quality vs. tracking accuracy
\end{itemize}

\section{Related Work}

\textbf{Physics-Based Character Control.} DeepMimic~\citep{peng2018deepmimic} pioneered using deep RL for physics-based character animation by combining imitation rewards with task rewards. AMP~\citep{peng2021amp} introduced adversarial motion priors, using a discriminator to distinguish policy motion from reference motion without explicit pose tracking rewards. ASE~\citep{peng2022ase} extended this to learn reusable skill embeddings from large motion datasets.

\textbf{Differential Discriminators.} ADD~\citep{zhang2025add} proposed using differential observations (relative differences between policy and reference) rather than absolute states for the discriminator. This approach enables the policy to track motion patterns while adapting to external forces or different morphologies.

\textbf{Robustness in Motion Imitation.} While several works address domain randomization for sim-to-real transfer~\citep{peng2018sim}, systematic studies of external force robustness in adversarial imitation learning remain limited. Our work fills this gap by evaluating how discriminator formulations affect robustness to sustained external forces.

\section{Methods}

\subsection{Background: AMP and ADD}

Both AMP and ADD use adversarial imitation learning where a discriminator $D$ tries to distinguish between policy-generated motions and reference motions, while the policy learns to fool the discriminator. The discriminator is trained via adversarial loss while the policy is trained via reinforcement learning with discriminator-based rewards.

\textbf{AMP}~\citep{peng2021amp} uses absolute state observations for discrimination. The discriminator $D_\theta: \mathcal{S} \rightarrow [0,1]$ is trained to distinguish policy states from reference states using binary cross-entropy loss. The style reward is computed as:
\begin{equation}
r_{style}(s_t) = -\log(1 - D_\theta(s_t))
\label{eq:amp_reward}
\end{equation}
where $s_t$ contains joint positions, velocities, and root information from the policy at time $t$. This reward encourages the policy to generate states that fool the discriminator (achieving $D_\theta(s_t) \rightarrow 1$).

\textbf{ADD}~\citep{zhang2025add} uses differential observations computed as the difference between policy and reference states. The discriminator observes:
\begin{equation}
\mathbf{o}^{diff}_t = \phi(s^{\pi}_t, s^{ref}_t)
\label{eq:add_obs}
\end{equation}
where $\phi$ computes relative features including root position difference $\Delta \mathbf{p}_{root}$, root velocity difference $\Delta \mathbf{v}_{root}$, and relative body positions with respect to the root. The discriminator $D_\theta: \mathcal{O}^{diff} \rightarrow [0,1]$ is trained with the zero vector $\mathbf{0}$ as the only positive sample (representing perfect tracking) and error vectors $\mathbf{o}^{diff}_t$ from the policy as negative samples. The reward encourages minimizing the differential:
\begin{equation}
r_{style}(s^{\pi}_t, s^{ref}_t) = -\log(1 - D_\theta(\mathbf{o}^{diff}_t))
\label{eq:add_reward}
\end{equation}
This differential formulation provides translation and rotation invariance, allowing policies to adapt their global position while maintaining relative motion patterns.

\subsection{Force Perturbation Framework}

We extended the MimicKit codebase to support configurable force perturbations:

\textbf{Random Force Model.} Forces are applied probabilistically with decay:
\begin{align}
\mathbf{f}_t &= \mathbf{f}_{t-1} \cdot \alpha_{decay} + \mathbf{f}_{new} \cdot \mathbb{I}(u < p_{force}) \label{eq:force_decay}\\
\mathbf{f}_{new} &\sim \mathcal{U}(-1, 1)^3 \odot \mathbf{s}_{force} \label{eq:force_sample}
\end{align}
where $\alpha_{decay} \in [0,1]$ is the decay factor (we use 0.95), $p_{force}$ is application probability (0.2 per step), $\mathbf{s}_{force}$ is the force scale vector, and $u \sim \mathcal{U}(0,1)$. The indicator function $\mathbb{I}(\cdot)$ ensures forces are applied stochastically.

\textbf{Curriculum Learning.} We gradually increase force magnitudes during training:
\begin{equation}
\mathbf{s}_{force}(t) = \min(\mathbf{s}_{max}, \mathbf{s}_{init} + \beta \cdot t)
\label{eq:curriculum}
\end{equation}
where $\beta$ controls the rate of increase, $t$ is the training iteration, and $\mathbf{s}_{max}$ is the target maximum force (e.g., [0, 0, 30N] or [10N, 10N, 30N]).

\subsection{Box-Carrying Dataset Generation}

We created reference motions for box-carrying by:
\begin{enumerate}
\item Retargeting human walking motion from AMASS~\citep{mahmood2019amass} to Unitree G1
\item Locking shoulder, elbow, and wrist joints to a fixed box-holding pose (arms extended forward)
\item Applying constant downward forces (0-100N) on wrist links during training to simulate box weight
\end{enumerate}

This generates naturalistic reference data where the lower body exhibits natural walking motion while the upper body maintains a carrying pose. The locked upper-body joints ensure the reference motion represents realistic box-carrying behavior.

\subsection{Steering Task with ADD}

Given AMP's position drift issues under force (Section~\ref{sec:results}), we developed a steerable policy using ADD to enable directional control while maintaining motion style. We implemented a steering environment where the robot must follow dynamically changing target directions while maintaining natural locomotion patterns.

\textbf{Dataset Preparation.} We retargeted the LAFAN dataset~\citep{harvey2020lafan} to the Unitree G1, using 6 diverse motion sets with neutral (natural) upper-body positions. This provides varied locomotion patterns including walking, running, and transitions.

\textbf{Discriminator Enhancement.} We augmented ADD's differential discriminator observations with velocity information, enabling the discriminator to assess motion quality in the context of current movement direction.

\textbf{Task Reward.} The steering task reward combines directional objectives:
\begin{equation}
r_{task} = w_{tar} \cdot r_{vel} + w_{face} \cdot r_{face}
\label{eq:steering_task}
\end{equation}
where $r_{vel} = \exp(-\|\mathbf{v} - v_{tar}\mathbf{d}_{tar}\|^2)$ rewards moving at target speed $v_{tar}$ in target direction $\mathbf{d}_{tar}$, and $r_{face} = \exp(-\|\mathbf{d}_{heading} - \mathbf{d}_{tar}\|^2)$ rewards facing the target direction. The total reward combines task and style objectives:
\begin{equation}
r_{total} = r_{task} + w_{style} \cdot r_{style}
\label{eq:total_reward}
\end{equation}

The choice of ADD over AMP for steering was motivated by ADD's superior position tracking under force, making it more suitable for directional control tasks despite overcompensation tendencies in single-axis force training.

\section{Experimental Setup}

\textbf{Environment.} We use Isaac Gym with the Unitree G1 humanoid (29 DoF) on flat terrain. Control frequency is 30Hz with simulation at 120Hz. Episode length is 10 seconds with early termination on falls.

\textbf{Training Details.} We use PPO~\citep{schulman2017ppo} with 4096 parallel environments. For AMP, we applied gradient clipping (max norm 1.0) and reduced discriminator strength (loss weight 0.5, gradient penalty 0.5) to prevent discriminator dominance. Training runs for 20,000-25,000 iterations (2-3 billion samples, 15-20 GPU hours on RTX 4090).

\textbf{Force Configurations.} We test four force scenarios:
\begin{itemize}
\item \textbf{Constant Light:} 30N downward force on wrists (simulating 3kg box)
\item \textbf{Constant Heavy:} 100N downward force on wrists (simulating 10kg box)
\item \textbf{Random + Constant:} 0-100N random forces with 0.95 decay plus constant weight
\item \textbf{Multi-Directional:} [10N, 10N, 30N] forces in x, y, z axes to reduce single-axis overcompensation
\end{itemize}

\textbf{Evaluation Metrics.} We measure:
\begin{itemize}
\item Episode length (surrogate for stability and training progress)
\item \textbf{Pose tracking error}: Mean joint-wise position error from reference motion
\item Root tracking error: Position and velocity deviation from reference trajectory
\item End-effector position error (feet and hands)
\item Discriminator reward (style matching quality)
\end{itemize}

\section{Results}
\label{sec:results}

We present our results in three stages: (1) baseline training dynamics with single-axis forces, (2) qualitative analysis revealing the overcompensation problem, and (3) quantitative comparison demonstrating that multi-directional forces solve this issue.

\subsection{Training Dynamics: Baseline [0,0,30N] Curriculum}

We first establish baseline performance by training both methods with a [0,0,30N] (downward-only) force curriculum.

\textbf{ADD Performance.} Figure~\ref{fig:add_results_part1} shows ADD training curves. Key observations:
\begin{itemize}
\item Episode length saturates quickly (within 5000 iterations), reaching maximum rollout duration
\item Discriminator maintains high loss, suggesting policy motion differs from reference despite achieving stability
\item Policy learns rapid force compensation, enabling continued walking
\end{itemize}

\textbf{AMP Performance.} Figure~\ref{fig:amp_results_part1} shows AMP training after hyperparameter tuning (gradient clipping, discriminator weakening):
\begin{itemize}
\item More gradual learning curve than ADD, reaching stability around 15,000 iterations
\item Discriminator reward steadily increases, indicating better style matching
\item Higher pose tracking accuracy but lower stability under force
\end{itemize}

\begin{figure}[t]
    \centering
    \subfigure[ADD episode length]{
        \includegraphics[width=0.3\linewidth]{data/ADD/add_episode_length.png}
        \label{fig:add_episode_length}
    }
    \subfigure[ADD discriminator]{
        \includegraphics[width=0.3\linewidth]{data/ADD/add_discmetrics.png}
    }
    \subfigure[ADD losses]{
        \includegraphics[width=0.3\linewidth]{data/ADD/add_losscurves.png}
    }
    \caption{\textbf{ADD training curves under 30N constant force.} Episode length saturates quickly (a), discriminator loss remains high (b), but policy converges (c).}
    \label{fig:add_results_part1}
\end{figure}

\begin{figure}[t]
    \centering
    \subfigure[AMP episode length]{
        \includegraphics[width=0.3\textwidth]{data/AMP/figure1_episode_length.png}
        \label{fig:amp_episode_length}
    }
    \subfigure[AMP discriminator]{
        \includegraphics[width=0.3\textwidth]{data/AMP/figure2_discriminator_metrics.png}
    }
    \subfigure[AMP losses]{
        \includegraphics[width=0.3\textwidth]{data/AMP/figure3_loss_curves.png}
    }
    \caption{\textbf{AMP training curves under force perturbations.} More gradual learning (a), improving discriminator reward (b), and stable loss convergence (c).}
    \label{fig:amp_results_part1}
\end{figure}

\subsection{Qualitative Analysis: The Overcompensation Problem}

Having established baseline training, we now investigate policy behavior under different force magnitudes to understand robustness limits.

\paragraph{Single-Axis Forces Cause Overcompensation in ADD}
Under [0,0,30N] curriculum training, ADD produces stable walking but develops unnatural compensation strategies:
\begin{itemize}
\item Arms raised above head to counterbalance downward forces
\item Exaggerated leaning motions during gait
\item \textbf{Persistent overcompensation}: Even at zero force, policy maintains defensive posture
\end{itemize}

The [0,0,100N] curriculum exacerbates this problem. Figure~\ref{fig:add_results} shows that the aggressive curriculum creates policies that overcompensate severely even without applied forces, sometimes becoming unstable. This behavior suggests the policy learns axis-specific compensation strategies (raising arms for downward forces) that persist inappropriately.

\paragraph{Multi-Directional Forces Solve Overcompensation}
To address this issue, we trained ADD with [10N, 10N, 30N] multi-directional forces. This configuration applies forces in x, y, and z axes, preventing axis-specific compensation strategies. Figure~\ref{fig:multi_directional_comparison} shows that multi-directional training produces natural poses even after force removal, validating our hypothesis that force diversity prevents overcompensation.

% TODO: Add Figure showing [0,0,30], [0,0,100], and [10,10,30] side-by-side

\begin{figure}[t]
    \centering
    \includegraphics[width=0.95\linewidth]{data/ADD/DRLmidpointsim_drawio.png}
    \caption{\textbf{ADD behavior under single-axis force curricula.} 30N curriculum allows adaptation but shows arm raising, while 100N curriculum causes severe overcompensation even at zero force.}
    \label{fig:add_results}
\end{figure}

\begin{figure}[t]
    \centering
    % TODO: Replace with actual multi-directional comparison figure
    % Should show 3 panels: [0,0,100] single-axis (overcompensation) vs [10,10,30] multi-directional (natural)
    \includegraphics[width=0.95\linewidth]{data/ADD/multi_directional_comparison.png}
    \caption{\textbf{Effect of multi-directional forces on ADD.} Policies trained with single-axis [0,0,100N] forces (left) overcompensate with raised arms at 0N, while multi-directional [10N,10N,30N] training (right) produces natural poses.}
    \label{fig:multi_directional_comparison}
\end{figure}

\paragraph{AMP: Better Pose Tracking, Lower Force Robustness}
AMP maintains more natural-looking poses and better pose tracking accuracy but exhibits lower force robustness. We trained AMP with the same three force configurations as ADD to enable direct comparison:

\textbf{Single-axis forces [0,0,30N] and [0,0,100N]:}
\begin{itemize}
\item Significant drift in root position during walking under force
\item Occasional stumbling and falls under moderate forces (>40N)
\item Better recovery when forces are removed compared to ADD
\item Maintains closer match to reference motion (higher discriminator rewards)
\end{itemize}

\textbf{Multi-directional forces [10,10,30N]:}
\begin{itemize}
\item Improved stability compared to single-axis training
\item Still exhibits position drift, making it less suitable for directional control
\item Maintains natural pose quality across all force magnitudes
\end{itemize}

Figure~\ref{fig:amp_results} shows that AMP can walk with light constant forces but experiences drift and instability with strong perturbations, confirming the trade-off between pose tracking accuracy and force robustness across all training configurations.

\begin{figure}[t]
    \centering
    \includegraphics[width=0.95\linewidth]{data/AMP/amp.png}
    \caption{\textbf{AMP behavior under constant forces.} Model trained with curriculum up to 100N shows reasonable walking with light forces but drifts and falls under strong forces.}
    \label{fig:amp_results}
\end{figure}

\subsection{Quantitative Comparison}

We now quantify the performance differences across force configurations and methods.

\paragraph{Pose Tracking and Stability Metrics}
Table~\ref{tab:pose_tracking} presents pose tracking error and episode length across force magnitudes for both AMP and ADD trained with all three force configurations. Key findings:
\begin{itemize}
\item \textbf{AMP}: Excellent pose tracking at 0N across all configs, but rapid degradation under force; multi-directional training provides modest improvement
\item \textbf{ADD}: Lower pose tracking error under force but higher error at 0N with single-axis training (overcompensation)
\item \textbf{Multi-directional benefit}: [10,10,30] training reduces ADD's 0N pose error while maintaining robustness; AMP shows smaller improvement but still drifts
\end{itemize}

\begin{table}[h]
\centering
\caption{Pose tracking error (mean $\pm$ std, degrees) and episode length at different force magnitudes}
\label{tab:pose_tracking}
\begin{tabular}{lcccc}
\toprule
\textbf{Method} & \textbf{0N} & \textbf{30N} & \textbf{50N} & \textbf{100N} \\
\midrule
\multicolumn{5}{l}{\textit{Pose Tracking Error (degrees)}} \\
AMP [0,0,30] & -- $\pm$ -- & -- $\pm$ -- & -- $\pm$ -- & -- \\
AMP [0,0,100] & -- $\pm$ -- & -- $\pm$ -- & -- $\pm$ -- & -- $\pm$ -- \\
AMP [10,10,30] & -- $\pm$ -- & -- $\pm$ -- & -- $\pm$ -- & -- \\
ADD [0,0,30] & -- $\pm$ -- & -- $\pm$ -- & -- $\pm$ -- & -- \\
ADD [0,0,100] & -- $\pm$ -- & -- $\pm$ -- & -- $\pm$ -- & -- $\pm$ -- \\
ADD [10,10,30] & -- $\pm$ -- & -- $\pm$ -- & -- $\pm$ -- & -- \\
\midrule
\multicolumn{5}{l}{\textit{Episode Length (seconds)}} \\
AMP [0,0,30] & -- $\pm$ -- & -- $\pm$ -- & -- $\pm$ -- & -- \\
AMP [0,0,100] & -- $\pm$ -- & -- $\pm$ -- & -- $\pm$ -- & -- $\pm$ -- \\
AMP [10,10,30] & -- $\pm$ -- & -- $\pm$ -- & -- $\pm$ -- & -- \\
ADD [0,0,30] & -- $\pm$ -- & -- $\pm$ -- & -- $\pm$ -- & -- \\
ADD [0,0,100] & -- $\pm$ -- & -- $\pm$ -- & -- $\pm$ -- & -- $\pm$ -- \\
ADD [10,10,30] & -- $\pm$ -- & -- $\pm$ -- & -- $\pm$ -- & -- \\
\bottomrule
\end{tabular}
\end{table}

% TODO: Fill in table with your experimental data for both AMP and ADD across all force configs

\paragraph{Method Comparison}
Table~\ref{tab:comparison} summarizes the key differences between AMP and ADD under force perturbations. Both methods were also evaluated on steering tasks, where they demonstrated the ability to follow target directions while maintaining learned motion styles, with ADD showing faster adaptation to directional control.

\begin{table}[t]
\centering
\caption{Comparison of AMP and ADD under force perturbations}
\label{tab:comparison}
\begin{tabular}{lcc}
\toprule
\textbf{Metric} & \textbf{AMP} & \textbf{ADD (single-axis)} \\
\midrule
Training speed & Slower (15k iter) & Faster (5k iter) \\
Motion naturalness & Better & Worse (overcompensation) \\
Pose tracking @0N & Excellent (low error) & Poor (overcompensation) \\
Pose tracking @force & Poor (drifts) & Better (maintains pattern) \\
Stability under force & Moderate (30-50N) & Good (50-100N) \\
Recovery after force & Good & Poor (persistent) \\
\midrule
\multicolumn{3}{l}{\textbf{ADD (multi-directional [10,10,30])}} \\
Motion naturalness & \multicolumn{2}{c}{Improved - natural poses at 0N} \\
Pose tracking @0N & \multicolumn{2}{c}{Improved - reduced overcompensation} \\
Stability under force & \multicolumn{2}{c}{Maintained - robust to varied forces} \\
\bottomrule
\end{tabular}
\end{table}

\section{Discussion}

\subsection{Why ADD Overcompensates (But Maintains Robustness)}

ADD's differential discriminator observes $\Delta s = s_\pi - s_{ref}$, which includes differences in body positions and velocities. When single-axis forces are applied:

\begin{enumerate}
\item The policy learns any pose that maintains similar \emph{relative} motion patterns to the reference, sacrificing absolute pose tracking accuracy for stability
\item Raised arms or exaggerated leans maintain differential consistency while counterbalancing forces, leading to poor pose tracking at 0N
\item The curriculum creates expectation of forces, causing persistent compensation strategies even when unnecessary
\end{enumerate}

This explains why ADD shows higher pose tracking error at 0N (overcompensation) but better error under force (robust patterns). Multi-directional forces [10,10,30] prevent axis-specific strategies, improving pose tracking at 0N while maintaining robustness.

\subsection{Why AMP Struggles with Pose Tracking Under Force}

AMP's absolute discriminator observes full state $s$, requiring close pose tracking to reference for high rewards. Under external forces:

\begin{enumerate}
\item The policy must balance discriminator reward (match reference pose) with task reward (avoid falling under force)
\item Forces cause accumulating root position errors that degrade overall pose tracking accuracy
\item Gradient clipping and discriminator weakening (necessary for training stability) reduce corrective signals for pose tracking
\end{enumerate}

The trade-off is excellent pose tracking at 0N but rapid degradation under force. AMP maintains natural motion quality but cannot sustain pose tracking accuracy under sustained perturbations.

\subsection{Curriculum Learning Insights}

Our force curriculum reveals important insights about force magnitude and directionality:

\textbf{Force Magnitude Effects:}
\begin{itemize}
\item \textbf{30N curriculum:} Produces policies that handle 0-40N robustly without overcompensation
\item \textbf{100N curriculum:} Creates defensive policies that overcompensate even at zero force
\end{itemize}

\textbf{Multi-Directional Forces:} To address the overcompensation problem observed in single-axis (downward-only) training, we experimented with multi-directional forces [10N, 10N, 30N] applied in x, y, and z axes. This configuration:
\begin{itemize}
\item Prevents the policy from learning axis-specific compensation strategies (e.g., raising arms for downward forces)
\item Encourages more balanced, natural responses to perturbations
\item Reduces persistent overcompensation when forces are absent
\item Results in better generalization across different force directions
\end{itemize}

This suggests that \emph{force diversity} during training is as important as force magnitude. Policies trained with multi-directional perturbations develop more robust and natural-looking compensation strategies compared to single-axis training.

\subsection{Limitations and Future Work}

\textbf{Current Limitations:}
\begin{itemize}
\item Evaluation limited to downward wrist forces; real box-carrying involves distributed forces
\item Flat terrain only; uneven ground would compound challenges
\item No real hardware validation; sim-to-real gap remains unquantified
\item Limited analysis of failure modes (what specific forces cause falls?)
\end{itemize}

\textbf{Future Directions:}
\begin{itemize}
\item \textbf{Hybrid discriminators:} Combine ADD's robustness with AMP's quality via multi-scale discrimination
\item \textbf{Force-aware curricula:} Include force magnitude as context input, with curriculum that increases then decreases forces
\item \textbf{Explicit pose constraints:} Add soft constraints on unnatural joint configurations
\item \textbf{Multi-task training:} Jointly train on steering, force perturbations, and terrain variations
\item \textbf{Real robot deployment:} Evaluate on physical Unitree G1 with actual box-carrying tasks
\end{itemize}

\section{Conclusion}

We presented a systematic study of AMP and ADD under external force perturbations, motivated by practical tasks like box-carrying. Through progressive experimentation with [0,0,30N], [0,0,100N], and [10,10,30N] force configurations, we uncovered fundamental trade-offs between these approaches and developed effective mitigation strategies.

Our key findings include:
\begin{enumerate}
\item \textbf{Discriminator formulation determines trade-offs}: ADD (differential) enables force robustness but sacrifices pose tracking accuracy at 0N through overcompensation; AMP (absolute) maintains excellent pose tracking at 0N but rapidly degrades under force
\item \textbf{Curriculum magnitude matters}: [0,0,100N] curricula induce persistent overcompensation, while [0,0,30N] provides better generalization
\item \textbf{Multi-directional forces solve overcompensation}: [10N,10N,30N] forces prevent axis-specific compensation strategies, improving pose tracking at 0N while maintaining robustness under force
\item \textbf{Method selection depends on task}: AMP's position drift makes it unsuitable for directional control tasks; ADD's robustness enables successful steering integration despite overcompensation in single-axis training
\item \textbf{Force diversity trumps magnitude}: Training with varied force directions produces more natural and generalizable policies than training with higher magnitude single-axis forces
\end{enumerate}

These insights inform the design of robust motion imitation systems for real-world humanoid robotics, where both naturalness and reliability under perturbations are essential. The extended MimicKit framework and evaluation protocols developed in this work provide a foundation for future research on robust physics-based character control.

\section*{Acknowledgments}

We thank the MimicKit authors for releasing their codebase and the course staff for valuable feedback. Computational resources were provided by Carnegie Mellon University.

\bibliographystyle{plain}
\begin{thebibliography}{10}

\bibitem{peng2018deepmimic}
Xue~Bin Peng, Pieter Abbeel, Sergey Levine, and Michiel van~de Panne.
\newblock Deepmimic: Example-guided deep reinforcement learning of
  physics-based character skills.
\newblock {\em ACM Transactions on Graphics}, 37(4):1--14, 2018.

\bibitem{peng2021amp}
Xue~Bin Peng, Ze~Ma, Pieter Abbeel, Sergey Levine, and Angjoo Kanazawa.
\newblock Amp: Adversarial motion priors for stylized physics-based character
  control.
\newblock {\em ACM Transactions on Graphics}, 40(4):1--15, 2021.

\bibitem{peng2022ase}
Xue~Bin Peng, Yunrong Guo, Lina Halper, Sergey Levine, and Sanja Fidler.
\newblock Ase: Large-scale reusable adversarial skill embeddings for physically
  simulated characters.
\newblock {\em ACM Transactions on Graphics}, 41(4):1--16, 2022.

\bibitem{zhang2025add}
Ziyu Zhang, Sergey Bashkirov, Dun~Yang, Yi~Shi, Michael Taylor, and Xue~Bin
  Peng.
\newblock Physics-based motion imitation with adversarial differential
  discriminators.
\newblock In {\em SIGGRAPH Asia 2025 Conference Papers}, 2025.

\bibitem{schulman2017ppo}
John Schulman, Filip Wolski, Prafulla Dhariwal, Alec Radford, and Oleg Klimov.
\newblock Proximal policy optimization algorithms.
\newblock {\em arXiv preprint arXiv:1707.06347}, 2017.

\bibitem{peng2018sim}
Xue~Bin Peng, Marcin Andrychowicz, Wojciech Zaremba, and Pieter Abbeel.
\newblock Sim-to-real transfer of robotic control with dynamics randomization.
\newblock In {\em IEEE International Conference on Robotics and Automation},
  pages 3803--3810, 2018.

\bibitem{mahmood2019amass}
Naureen Mahmood, Nima Ghorbani, Nikolaus~F. Troje, Gerard Pons-Moll, and
  Michael~J. Black.
\newblock {AMASS}: Archive of motion capture as surface shapes.
\newblock In {\em IEEE International Conference on Computer Vision (ICCV)},
  pages 5442--5451, October 2019.

\bibitem{harvey2020lafan}
F{\'e}lix~G. Harvey, Mike Yurick, Derek Nowrouzezahrai, and Christopher Pal.
\newblock Robust motion in-betweening.
\newblock {\em ACM Transactions on Graphics}, 39(4), 2020.

\end{thebibliography}

\end{document}
